\documentclass{article}
\usepackage{graphicx}
\usepackage{amsmath}
\begin{document}
\title{A Realistic Generic Article on Semantic Conversion}
\author{Example Author}
\maketitle
\begin{abstract}
This fixture exercises a generic article with sections, equations, tables, figures, and references.
\end{abstract}
\section{Introduction}
Generic manuscripts often mix prose, mathematical notation, and cross references. Figure~\ref{fig:generic-pipeline} and Equation~\ref{eq:generic-score} provide representative anchors.
\section{Method}
\begin{equation}
  s(d) = \alpha q(d) + (1-\alpha)c(d) \label{eq:generic-score}
\end{equation}
\begin{table}[ht]
  \centering
  \begin{tabular}{lrr}
  Method & Score & Time\\
  Baseline & 71.2 & 3.4\\
  Proposed & 84.6 & 2.8\\
  \end{tabular}
  \caption{Generic evaluation summary.}
\end{table}
\begin{figure}[ht]
  \centering
  \includegraphics[width=.5\textwidth]{generic-pipeline.png}
  \caption{Generic conversion pipeline.}
  \label{fig:generic-pipeline}
\end{figure}
\section{Conclusion}
This document is intentionally compact but closer to a real article than the minimal smoke fixture.
\bibliographystyle{plain}
\bibliography{refs}
\end{document}
